\section{Annihilators and the closed-fiber representation}\label{homol}
Let $R\subseteq S$ be a module-finite extension of a local ring with $\pd_R S\le1$. Fix the minimal resolution and notation
\begin{equation}\label{homol-resolution}
 0\longrightarrow R^b\xrightarrow{d}R^a
   \xrightarrow{\pi}S\longrightarrow0,
 \qquad H=\Ext^1_R(S,R),\qquad F=S/\m S.
\end{equation}
At each minimal prime of $R$, the localization of $S$ is free by
Auslander--Buchsbaum, since the localized base has depth zero.  Its rank
is $r=a-b$, and injectivity of $R\to S$ gives $r>0$.  Minimality therefore
gives
\[
 \dim_k(H/\m H)=b<a=\dim_k F.
\]
If $b=0$, the module $S$ is free and $1_S\notin\m S$ by Nakayama's lemma, so a coordinate functional gives a retraction.
We may consequently assume $b>0$ whenever we argue by contradiction.
Put $V=H/\m H$ and write $D_k(-)=\Hom_k(-,k)$.  Multiplication maps on
$S$ act on $H$ by contravariance. Commutativity makes this an
$S$-module action.  The induced actions on $V$ and on Tor factor
through $F$.

\begin{lemma}[The Tor obstruction]\label{torob}
If $R\subseteq S$ does not split, then
\[
 \operatorname{Tor}_1^R(H,k)\simeq D_kF
\]
as $F$-modules, and $V$ is not an $F$-module direct summand of this
Tor group.  Moreover,
\[
 \operatorname{Ann}_S H\subseteq J S,
 \qquad J=\{\phi(1_S):\phi\in\Hom_R(S,R)\}.
\]
\end{lemma}
\begin{proof}
Nonsplitting means $J\subseteq\m$.  In fact every $R$-linear functional
$\phi:S\to R$ then takes all its values in $\m$: for $s\in S$, the
functional $t\mapsto\phi(st)$ takes the value $\phi(s)$ at $1_S$.
Dualizing \eqref{homol-resolution} gives
\[
 0\longrightarrow S^*\longrightarrow(R^a)^*
   \xrightarrow{d^*}(R^b)^*\longrightarrow H\longrightarrow0.
\]
The entries of $d^*$ lie in $\m$, and the preceding observation gives
$S^*\subseteq\m(R^a)^*$.  Thus the two displayed free modules are the
first two modules of a minimal free resolution of $H$.  It follows that
\[
 \operatorname{Tor}_1^R(H,k)=(R^a)^*\otimes_R k
    \simeq D_k(S/\m S).
\]
A lift of multiplication by $s\in S$ to \eqref{homol-resolution}
acts on this dual presentation by transposition.  The identification
is therefore $F$-linear.

Suppose now that $V$ were an $F$-module summand of $D_kF$.
Duality would make $D_kV$ a summand of the regular $F$-module $F$,
so $D_kV\simeq eF$ for an idempotent $e\in F$.  In particular $e$
acts as the identity on both $D_kV$ and $V$.  Choose a lift $s\in S$
of $e$.  Multiplication by $s$ on $H$ is surjective by Nakayama's
lemma, because it induces the identity on $V$.  A surjective
endomorphism of a Noetherian module is an automorphism.  It follows
that $s$ acts invertibly on $\operatorname{Tor}_1^R(H,k)$.
Under the preceding identification this action is the action of $e$
on $D_kF$.  That action is faithful, and an invertible idempotent
endomorphism is the identity.  Hence $e=1$ in $F$, which would give
$\dim_kV=\dim_kF$, contrary to $b<a$.

Finally, let $s\in\operatorname{Ann}_S H$, and lift multiplication
by $s$ to matrices $P$ on $R^a$ and $Q$ on $R^b$, with
$Pd=dQ$.  Since $Q^*$ induces zero on $\operatorname{coker}d^*$,
its image is contained in $\operatorname{im}d^*$.  Freeness of
$(R^b)^*$ permits a lift $W:(R^b)^*\to(R^a)^*$ such that
$Q^*=d^*W$.  Setting $L=W^*$ gives $Q=Ld$.  Consequently
$(P-dL)d=0$, and there is a map $g:S\to R^a$ with
$P-dL=g\pi$.  Applying $\pi$ shows that $\pi g$ is multiplication
by $s$.  If $g_1,\ldots,g_a$ are the coordinate functionals of $g$,
then
\[
 s=\pi g(1_S)=\sum_{i=1}^a g_i(1_S)\pi(e_i)\in JS,
\]
as required.
\end{proof}

\subsection{An annihilator with nonzero linear term}
\begin{theorem}\label{linear-ann}
If $\operatorname{Ann}_R H\not\subseteq\m^2$, the extension splits.
On a one-dimensional Gorenstein base it suffices that $\m H^0_{\m}(S)=0$.
\end{theorem}
\begin{proof}
If $H=0$, minimality of the dual presentation forces $b=0$, and the
extension splits. Otherwise every annihilator of $H$ is a nonunit.
The hypothesis therefore supplies
$x\in\operatorname{Ann}_R H\cap(\m\setminus\m^2)$.
Choose a minimal free cover $P\to H$, with kernel $K$.  Then
\[
 xP\subseteq K\subseteq\m P,
 \qquad V=P/\m P,
 \qquad \operatorname{Tor}_1^R(H,k)=K/\m K.
\]
Multiplication by $x$ induces a well-defined map
\[
 \iota_x:V\longrightarrow K/\m K,
 \qquad \bar v\longmapsto xv\bmod\m K,
\]
because $x\m P\subseteq\m K$.  Choose a $k$-linear functional
$\lambda:\m/\m^2\to k$ with $\lambda(\bar x)=1$.
The inclusion $K\subseteq\m P$ gives a map
\[
 K/\m K\longrightarrow\m P/\m^2P
       \simeq(\m/\m^2)\otimes_k V.
\]
Composing it with $\lambda\otimes1_V$ gives a left inverse of
$\iota_x$.

Both maps are $F$-linear.  Indeed a lift to $P$ of multiplication
by an element of $S$ preserves $K$ and commutes with the scalar $x$. On $\m P/\m^2P$ it acts as the identity on the first tensor factor
and by its action on $V$ on the second.  Thus $V$ would be a direct
summand of $\operatorname{Tor}_1^R(H,k)$ in the nonsplit case,
contradicting Lemma~\ref{torob}.

For the stated consequence on a one-dimensional Gorenstein base,
set
\[
 T=H^0_{\m}(S)
   =\{s\in S:\m^n s=0\text{ for some }n\geq0\}.
\]
This is a finite-length module.  After completion, local duality
identifies $\Ext^1_{\widehat R}(S\otimes_R\widehat R,\widehat R)$
with the Matlis dual of $T$.  Here the Matlis dual is
$\Hom_{\widehat R}(T,E_{\widehat R}(k))$, with $E_{\widehat R}(k)$
the injective hull of the residue field.  Matlis duality is faithful
on finite-length modules, so for every $h\geq1$
\begin{equation}\label{torsion-annihilator-equivalence}
 \m^hT=0\quad\Longleftrightarrow\quad\m^hH=0.
\end{equation}
Ext commutes with this flat base change, and faithful flatness gives
the assertion over $R$.  If $\m T=0$, any minimal generator of
$\m$ is an annihilator with nonzero linear term.  A retraction over
the completion descends because its evaluation ideal is
$J\widehat R$, so $J\widehat R=\widehat R$ implies $J=R$.
\end{proof}

\subsection{The size of the closed-fiber action}
The numerical criterion can be stated without a restriction on the
number of algebra generators of the closed fiber.  What matters is
the dimension of the representation through which that fiber acts.

\begin{proposition}\label{representation-bound}
In the notation of \eqref{homol-resolution}, assume $b>0$ and
$\m^hH=0$ for some $h\geq1$.  Put
\[
 c=\dim_k\operatorname{im}
       \bigl(F\longrightarrow\operatorname{End}_k(V)\bigr).
\]
If
\begin{equation}\label{representation-inequality}
 a-c>b\,\ell_R(\m/\m^h),
\end{equation}
then $R\subseteq S$ splits.
\end{proposition}
\begin{proof}
The kernel of the representation has dimension $z=a-c$.
Since $V\ne0$, this kernel does not contain $1_F$.  Choose minimal
module generators $1_S,s_1,\ldots,s_{a-1}$ of $S$ so that the
classes of $s_1,\ldots,s_z$ form a basis of that kernel.
Write the rows of the matrix $d$ as
$\alpha,w_1,\ldots,w_{a-1}$, in this order, and set
\[
 N=R^{1\times b}/\langle w_1,\ldots,w_{a-1}\rangle,
 \qquad \eta=[\alpha]\in N.
\]
A functional taking $1_S$ to $1$ exists exactly when $\alpha$ lies
in the span of the other rows, or equivalently when $\eta=0$.
The dual presentation also identifies $H=N/R\eta$.  Hence
\begin{equation}\label{numerical-row-annihilator}
 \m^hN\subseteq R\eta.
\end{equation}

The rows of $d$ have entries in $\m$.  Consider the images of
$w_1,\ldots,w_z$ in $(\m/\m^h)^b$.  Their span has length at most
$b\ell_R(\m/\m^h)$, which is strictly smaller than $z$ by
\eqref{representation-inequality}.  In particular these $z$
generators cannot be minimal.  There is a relation among them
with a unit coefficient.  Complete its coefficient row to an
invertible matrix over $R$ and change the first $z$ rows accordingly.
After making the inverse change in the corresponding generators
of $S$, we may assume
\[
 w_1\in\m^h R^{1\times b}.
\]
The new $s_1$ still acts as zero on $V$, since its residue class is
a linear combination of the same kernel basis.

Lift multiplication by this $s_1$ to a pair of matrices $P,Q$ with
$Pd=dQ$.  The column of $P$ indexed by $1_S$ can be chosen to be
the basis vector indexed by $s_1$.  On the dual presentation of
$H$, the matrix $Q^*$ induces multiplication by $s_1$. Since
$d^*$ is minimal, its reduction modulo $\m$ is exactly the action
on $V$.  Thus $Q\equiv0\pmod\m$.  The row of $Pd=dQ$ indexed by
$s_1$ now reads
\[
 \alpha+\sum_{j=1}^{a-1}P_{1j}w_j=w_1Q.
\]
The right side belongs to $\m^{h+1}R^{1\times b}$.  Passing to
$N$ and using \eqref{numerical-row-annihilator}, we obtain
\[
 \eta\in\m^{h+1}N\subseteq\m R\eta.
\]
Nakayama's lemma applied to the cyclic module $R\eta$ gives
$\eta=0$, proving splitting.
\end{proof}

\begin{corollary}\label{schur-bound}
With the notation of Proposition~\ref{representation-bound}, put
$r=a-b$ and $d_h=\ell_R(\m/\m^h)$. Each of the following conditions suffices for splitting.
\begin{enumerate}[label=\textup{(\roman*)}]
\item $r>b d_h+\lfloor b^2/4\rfloor+1-b$.
\item $r>b d_h$ and $F$ is generated by at most two elements as a $k$-algebra.
\end{enumerate}
For $b\le3$, the bound $r>b d_h$ suffices for every closed fiber.
\end{corollary}
\begin{proof}
The image of $F$ is a commutative subalgebra of $\operatorname{End}_k(V)$, so the Schur--Jacobson bound gives
$c\le\lfloor b^2/4\rfloor+1$ \cite{Jac}. For two generators Gerstenhaber's theorem gives $c\le b$ \cite{Gerst,Berg}.
Both bounds hold over $k$, since scalar extension preserves the dimension of the image. Substitute in \eqref{representation-inequality}. For $b=1,2,3$ the first upper bound equals $b$.
\end{proof}

\subsection{Lifting a retraction from the tangent cone}\label{tangent}
\begin{theorem}\label{tangent-theorem}
Suppose $\operatorname{gr}_{\m}R$ is a geometrically integral splinter over $k$. If
$(\operatorname{Jac} S)^2 H=0$, the extension splits.
This applies in particular when
$R=\mathbf Z_p[[u,v,z,w_1,\ldots,w_d]]/(uv-z^2+h)$,
where $p$ is any prime and $h\in(p,u,v,z,w_1,\ldots,w_d)^3$.
\end{theorem}
\begin{proof}
First assume that $S$ is local, with maximal ideal $\mathfrak n$
and residue field $l$, and suppose the extension does not split.
We prove
\begin{equation}\label{strong-annihilator}
 \operatorname{Ann}_S H\subseteq\mathfrak n\m S.
\end{equation}
Lemma~\ref{torob} first gives
$\operatorname{Ann}_S H\subseteq JS\subseteq\m S$.
Choose minimal generators $x_1,\ldots,x_v$ of $\m$, and write an
annihilator as $u=\sum_i x_i a_i$ with $a_i\in S$.
Let $P\to H$ be a minimal free cover with kernel $K$, as above,
and lift multiplication by $a_i$ to $P_i:P\to P$.
Since $uH=0$, the operator $U=\sum_i x_iP_i$ takes $P$ into $K$
and induces a map
\[
 \bar U:V\longrightarrow K/\m K.
\]
This map is $F$-linear.  Indeed, if $P_s$ lifts multiplication by
$s\in S$, then $[P_s,P_i]$ maps $P$ into $K$, since the induced
actions commute on $H$.  Therefore
$[P_s,U]P\subseteq\m K$.
Contracting $K/\m K\to(\m/\m^2)\otimes_kV$ against the coordinate
functional dual to $x_i$ gives an $F$-linear map $\rho_i$ with
\[
 \rho_i\bar U=a_i|_V.
\]
If $a_i$ were a unit of $S$, its action on $V$ would be invertible. Composing $\rho_i$ with its inverse would split off $V$ from
$K/\m K$.  Lemma~\ref{torob} forbids this.  Thus each
$a_i\in\mathfrak n$, proving \eqref{strong-annihilator}.
The hypothesis $\mathfrak n^2H=0$ now gives
\begin{equation}\label{minimal-target-filtration}
 \mathfrak n^2=\m\mathfrak n,
 \qquad \mathfrak n^{j+1}=\m^j\mathfrak n\quad(j\geq1).
\end{equation}
Here the reverse inclusion $\m\mathfrak n\subseteq\mathfrak n^2$
is automatic, and the higher equalities follow by induction.

Put $P'=\operatorname{gr}_{\m}R$ and
$G=\operatorname{gr}_{\mathfrak n}S$.
Equation~\eqref{minimal-target-filtration} shows that $G$ is generated
in degrees zero and one as a finite $P'$-module.  It is also finite
over the graded algebra $l\otimes_kP'$, through the natural map to
$G$.  Geometric integrality makes $l\otimes_kP'$ a domain.  Since
$R\subseteq S$ is integral and $l/k$ is finite,
\[
 \dim(l\otimes_kP')=\dim R=\dim S=\dim G.
\]
A nonzero ideal in the affine domain $l\otimes_kP'$ has quotient
of smaller dimension, whereas finiteness gives
$\dim G=\dim((l\otimes_kP')/\ker)$.  Hence
\begin{equation}\label{graded-residue-inclusion}
 l\otimes_kP'\hookrightarrow G.
\end{equation}
In particular $P'\hookrightarrow G$ is a module-finite extension.
The splinter hypothesis gives a $P'$-linear retraction $G\to P'$.
Taking its component of matching degree on each homogeneous piece
produces a graded retraction $\gamma:G\to P'$.

Write $\lambda:l\to k$ for the degree-zero part of $\gamma$;
then $\lambda(1)=1$.  Choose a $k$-basis of $l$ beginning with $1$
and with the remaining basis elements in $\ker\lambda$.
Lift it to $\sigma_0=1_S,\sigma_1,\ldots,\sigma_{t-1}\in S$.
Their classes in $S/\m S$ are linearly independent, and we may
complete them to a basis by elements $y_1,\ldots,y_q\in\mathfrak n$.
On the corresponding free cover define an $R$-linear functional by
\[
 f(\sigma_0)=1,\qquad f(\sigma_i)=0\ (i>0),
 \qquad f(y_j)=L_j\in\m,
\]
where $L_j$ lifts $\gamma([y_j]_1)\in\m/\m^2$.

Consider any relation
\[
 \sum_i r_i\sigma_i+\sum_j a_jy_j=0.
\]
All its coefficients are in $\m$, by minimality of the generators.
In degree one of $G$, the second sum vanishes, since
$\m\mathfrak n=\mathfrak n^2$.  By
\eqref{graded-residue-inclusion} and independence of the residue
classes of the $\sigma_i$, we obtain $r_i\in\m^2$ for every $i$.
Taking degree two of the relation and applying the $P'$-linear
map $\gamma$ gives
\[
 r_0+\sum_j a_jL_j\in\m^3.
\]
Equivalently, in the notation of \eqref{homol-resolution},
$d^*f\in\m^3(R^b)^*$.
But $\m^2H=0$, so
\[
 \m^3(R^b)^*\subseteq\m\operatorname{im}d^*
                      =d^*(\m(R^a)^*).
\]
Choose $v\in\m(R^a)^*$ with $d^*v=d^*f$.
Then $f-v$ descends to $S$ and takes a unit value at $1_S$.
Normalizing that value gives a retraction, a contradiction.

For an arbitrary finite target $S$, pass to the completion of $R$.
The tangent cone is unchanged.  Since it is a domain, the completed
base is a domain because the initial forms of two nonzero elements have
nonzero product, and their orders are finite by Krull intersection.
The completed algebra is a product of finite local algebras.
At least one factor map from the completed base is injective.
Otherwise the product of the nonzero kernels would be a nonzero
ideal of that domain in the kernel of the map to the full product.  The Jacobson radical commutes
with this completion because in either algebra it is the inverse image
of the radical of the same finite-dimensional closed fiber.
Flat base change for Ext and passage to a direct factor preserve
both the projective-dimension hypothesis and the radical-square
annihilation.  Apply the local case to an injective factor, compose
with projection from the product, and descend using the evaluation
ideal as above.

Finally, consider the mixed-characteristic family in the statement.
The associated graded ring of the ambient regular local ring is
the polynomial ring
$\mathbf F_p[P,U,V,Z,W_1,\ldots,W_d]$.
The defining equation has initial form $UV-Z^2$.  Orders of
nonzero elements add in the ambient associated graded domain,
so the initial ideal of this principal ideal is generated by
$UV-Z^2$.  Thus the tangent cone is a polynomial extension of
\[
 k[U,V,Z]/(UV-Z^2)\simeq k[s^2,t^2,st],
 \qquad k=\mathbf F_p.
\]
This description remains a domain after every field extension.
Projection of $k[s,t]$ onto the monomials of even total degree is
linear over $k[s^2,t^2,st]$ and gives a retraction even when
$p=2$.  The same projection extends after adjoining the remaining
polynomial variables.  Polynomial rings over $k$ are splinters by
the coefficient-extraction argument of Section~\ref{geometry}.
The direct summand is a splinter by Lemma~\ref{conductor-lower}.
Finally $p$ is a nonzerodivisor on the stated hypersurface.
Indeed the defining equation remains nonzero modulo $p$ in the
ambient power series domain, so reduction modulo $p$ shows
$(uv-z^2+h):p=(uv-z^2+h)$.
All the asserted hypotheses are therefore satisfied.
\end{proof}
